%%=============================================================================
%% Voorwoord
%%=============================================================================

\chapter*{\IfLanguageName{dutch}{Woord vooraf}{Preface}}%
\label{ch:voorwoord}

Deze bachelorproef komt voort uit mijn interesse in mainframebeveiliging en toegangsbeheer. Tijdens mijn opleiding merkte ik dat beveiligingsproblemen niet altijd starten met een grote fout. Soms gaat het om kleine rechten die blijven staan. Een gebruiker krijgt tijdelijk toegang, verandert later van taak, maar de oude groep blijft bestaan. Na een tijd is niet meer duidelijk waarom die toegang er nog is.

Daarom vond ik privilege creep een interessant onderwerp. Het is technisch, maar het heeft ook veel te maken met dagelijks beheer. Rechten worden meestal met een goede reden gegeven. Het probleem ontstaat wanneer niemand later controleert of die rechten nog nodig zijn.

RACF en z/OS waren voor mij geen eenvoudige onderwerpen. Het mainframedomein heeft veel eigen termen en gewoontes. Toch vond ik het net daarom nuttig om dit te onderzoeken. Mainframes spelen nog altijd een grote rol in banken, overheden en andere grote organisaties. Goed toegangsbeheer blijft daar dus belangrijk.

Met deze bachelorproef wilde ik geen oplossing bouwen die alles automatisch oplost. Dat zou niet realistisch zijn. Mijn doel was om te onderzoeken hoe RACF data kan helpen om opvallende groepscombinaties te vinden. De PoC moet beheerders en auditors helpen om sneller te zien welke gebruikers verder bekeken moeten worden.

Ik wil mijn promotor Leendert Blondeel bedanken voor de begeleiding en feedback. Ook mijn co-promotor Jeroen Souverijn wil ik bedanken voor de inhoudelijke hulp en het meedenken over RACF en mainframebeveiliging.

Daarnaast bedank ik iedereen die tijd maakte om uitleg te geven, feedback te geven of mee te denken. Veel inzicht in mainframebeveiliging komt niet alleen uit documentatie, maar ook uit gesprekken met mensen die zulke omgevingen kennen.

Tot slot wil ik mijn familie, vrienden en medestudenten bedanken voor de steun tijdens het schrijven van deze bachelorproef.
